% 06_conclusion.tex
%
% Purpose:
%   Close the reorganized report and identify the main remaining tasks.
%
% Main algorithm:
%   No algorithm is introduced.  The section summarizes the formulation and the
%   validation work still needed.
%
% Based on:
%   sections/conclusion.tex from the first report draft.
%
% Main changes:
%   Revises the conclusion to match the new six-section paper organization.
%
% Numerical goal:
%   Clarify what must be completed before the draft becomes a numerical paper.

\section{Conclusion}

This report has organized the current Transmission Eigenvalue / Bloch-mode /
line-defect formulation into a compact paper-style draft.  The main numerical
structure is a center-cell boundary integral equation coupled to compact port
matching conditions.  The periodic half-lead exterior condition is imposed by
requiring the port Cauchy data to lie in Bloch-selected incoming or outgoing
trace subspaces, rather than by explicitly forming a full half-guide
Dirichlet-to-Neumann matrix.

The formulation is practical and finite-dimensional.  It separates continuous
modeling assumptions from numerical approximations and identifies several
heuristic choices, especially Rayleigh truncation and mode selection, that need
convergence checks.

The main remaining tasks are to add numerical experiments, document unit-circle
mode selection in more detail, verify any unresolved bibliography metadata, and
insert convergence tables and field plots for representative forced-scattering
and homogeneous singular-value scans.
